% ============================================================
% UNIDAD 10: INSTALACIÓN, CONFIGURACIÓN Y DOCUMENTACIÓN
% ============================================================
\clearpage
\section*{UNIDAD 10: INSTALACIÓN DE UN PANEL DIGITAL}
\subsection*{Planificación, tablero, cableado, configuración y documentación}

\begin{recuadro}
\textbf{Objetivo de la unidad:} planificar y ejecutar la instalación de instrumentos digitales:
diseño del panel, planos de cableado, armado del mazo, verificaciones antes de energizar,
configuración, calibración y documentación de la alteración.
\end{recuadro}

% ------------------------------------------------------------
\section*{1. PLANIFICACIÓN}
\begin{enumerate}
  \item \textbf{Verificar la elegibilidad:} que la aeronave esté en la lista del STC y que el STC
        esté aceptado por la autoridad; o, en una experimental, los requisitos aplicables.
  \item \textbf{Reunir la documentación:} manual de instalación del equipo, documentos del STC,
        manual de mantenimiento y de cableado de la aeronave.
  \item \textbf{Análisis de carga eléctrica:} que el alternador y las barras soporten el nuevo
        consumo (y lo que se retira).
  \item \textbf{Peso y balanceo:} calcular el cambio de peso y de centro de gravedad.
  \item \textbf{Lista de materiales} con números de parte, y \textbf{plan de trabajo} con tiempos.
\end{enumerate}

% ------------------------------------------------------------
\section*{2. EL PANEL DE INSTRUMENTOS}
\begin{itemize}
  \item \textbf{Ubicación:} el PFD frente al piloto, dentro de su campo visual primario; el
        instrumento de reserva visible desde la posición del piloto (criterios de la AC 23.1311-1).
  \item \textbf{Recortes:} con plantillas del fabricante; en paneles estructurales, verificar que
        el recorte esté aprobado.
  \item \textbf{Ventilación:} espacio detrás del panel y ventiladores cuando el equipo los exige.
  \item \textbf{Iluminación:} los equipos digitales regulan su brillo automáticamente o desde el
        reóstato del panel; la iluminación de los instrumentos restantes y de las placas debe
        quedar uniforme.
  \item \textbf{Placas y rótulos:} interruptores, disyuntores y limitaciones nuevas, según el STC.
\end{itemize}

% ------------------------------------------------------------
\section*{3. EL PLANO DE CABLEADO}

A partir de los diagramas de interconexión del manual de instalación se dibuja el
\textbf{plano de cableado de la aeronave}: cada cable con su identificación, sección, origen
(conector y número de contacto), destino, blindaje y conexión a masa.

\begin{center}
\small
\begin{tabular}{|c|c|c|c|c|p{0.18\linewidth}|}
\hline
\textbf{Cable} & \textbf{AWG} & \textbf{Desde} & \textbf{Hasta} & \textbf{Blindaje} &
\textbf{Señal} \\ \hline
P1-01 & 22 & J1-1 (equipo) & CB ``PFD'' & -- & Alimentación \\ \hline
P1-02 & 22 & J1-2 & Masa GS1 & -- & Masa \\ \hline
P1-10 & 22 & J1-10 & GMU J1-3 & Sí, un extremo & Datos (RS-485 A) \\ \hline
\ldots & & & & & \\ \hline
\end{tabular}
\end{center}
\textit{(Ejemplo ilustrativo de tabla de cableado; los datos reales salen del manual de
instalación.)}

Herramientas útiles: software de dibujo eléctrico o de diseño de mazos, y planillas de cálculo
para las tablas de cables. El plano final queda en la documentación de la aeronave: es lo que
permitirá a otro técnico reparar el sistema dentro de diez años.

% ------------------------------------------------------------
\section*{4. ARMADO Y VERIFICACIÓN DEL MAZO}
\begin{enumerate}
  \item Armar el mazo en el tablero (Unidad 6), crimpar contactos y terminales (Unidad 7) y armar
        los conectores (Unidad 8).
  \item \textbf{Verificación pin a pin:} con el plano y un multímetro (o un probador de mazos),
        verificar cada conexión: continuidad de cada cable, ausencia de cortocircuitos entre
        cables y contra masa, blindajes conectados donde corresponde.
  \item Instalar el mazo en la aeronave con sus abrazaderas, pasacables y separaciones.
  \item \textbf{Verificación antes de conectar los equipos:} con los conectores de los equipos
        \textbf{desconectados}, energizar y medir que en cada pin de alimentación haya la tensión
        correcta y que no haya tensión donde no debe haberla. Luego desenergizar y conectar los
        equipos. Así se evita quemar un equipo por un error de cableado.
\end{enumerate}

% ------------------------------------------------------------
\section*{5. CONFIGURACIÓN Y CALIBRACIÓN}
\begin{itemize}
  \item \textbf{Software:} cargar la versión aprobada.
  \item \textbf{Configuración de la aeronave:} tipo de aeronave, velocidades y marcas de los arcos
        (según el manual de vuelo), unidades, interfaces (puertos serie, ARINC 429 y su velocidad,
        buses), equipos conectados.
  \item \textbf{Sensores de motor:} tipo de cada sensor, rangos, colores, alertas y calibración de
        combustible (llenando los tanques por etapas y registrando cada punto).
  \item \textbf{Calibración del AHRS:} orientación de montaje y nivelación (cabeceo y alabeo) con la
        aeronave en posición de nivelación.
  \item \textbf{Calibración del magnetómetro:} según el procedimiento del fabricante, generalmente
        girando la aeronave en una zona sin interferencias magnéticas.
  \item \textbf{Bases de datos:} cargar las vigentes.
  \item Guardar una copia de la configuración final.
\end{itemize}

% ------------------------------------------------------------
\section*{6. ENSAYOS Y DOCUMENTACIÓN FINAL}
\begin{itemize}
  \item Ensayos en tierra del manual de instalación: funcionamiento de cada función, interfaces
        con radios, GPS, transpondedor y piloto automático, alertas, modos reversionarios y
        batería de reserva.
  \item \textbf{Ensayo pitot-estático} y de altímetro/transpondedor (Unidad 11).
  \item \textbf{Ensayo de compatibilidad electromagnética:} con todos los sistemas funcionando,
        verificar que el nuevo equipo no interfiera con los existentes (y viceversa), por ejemplo
        transmitiendo con las radios y observando las pantallas y el magnetómetro.
  \item Ensayo en vuelo si el STC o el manual lo requiere.
  \item \textbf{Documentación:} registro de la alteración en los registros de mantenimiento,
        formulario de alteración mayor que exija la autoridad, actualización del peso y balanceo,
        \textbf{suplemento del manual de vuelo}, instrucciones de aeronavegabilidad continuada
        (ICA) y planos de cableado actualizados.
\end{itemize}

% ------------------------------------------------------------
\section*{7. ACTIVIDADES DE LA UNIDAD}
\begin{enumerate}
  \item Hacer el plan completo de instalación de un G5 como horizonte artificial en un Piper PA-28:
        elegibilidad, materiales, carga eléctrica, peso, tiempos.
  \item Dibujar el plano de cableado del G5 con su magnetómetro a partir del manual de instalación.
  \item Diseñar en escala un panel con un G3X Touch de 10 pulgadas, un G5 de reserva, radio y
        transpondedor.
  \item Redactar el registro de mantenimiento de la instalación.
\end{enumerate}
